\subsection{Task resolution, positive error, and the complete state}
\label{sec:v15-task}
We give a quantitative state law at a source-consuming restart. It
complements the endogenous zero-face theorem; it does not replace the
causal update constraint by a terminal partition.

Let $\mathcal H$ be a finite set of supported histories at a fixed
checkpoint. Choose $d$ actual marked continuation-event prediction
tasks, each including its declared open-loop continuation control word.
For $h\in\mathcal H$ let $v(h)\in[0,1]^d$ be the vector of their
conditional probabilities. A source-consuming statistic has rows
$Q(m\mid h)$, $m\in[K]$. For every preparation $\lambda$ on
$\mathcal H$, the task is announced after storage and its forecast is
scored by squared error against the actual event. Define
\[
 e(Q)=\sup_\lambda\max_{1\le j\le d}
       \E_\lambda\Var(v_j(H)\mid M),\qquad
 e_K=\inf_Q e(Q).
\]
The expression is precisely excess Brier risk over the full-history
forecast: conditional on $H$, event variance is an unavoidable common
term, and the law of total variance gives the displayed residual.
Decoders may use the announced preparation and task, but the same
statistic $Q$ must work for all of them. Let $N(\eta)$ be the maximum
size of a subset of $v(\mathcal H)$ whose Euclidean pair distances are
at least $\eta>0$. Let $C(r)$ be the minimum number of sets covering
$v(\mathcal H)$ with coordinatewise diameter at most $r$.

\begin{theorem}[Positive-error state bounds]\label{thm:v15-task}
For all positive $K$ and $\eta$,
\begin{equation}\label{eq:v15-task-packing}
 e_K\ge\frac{\eta^2}{2d}\left(1-\frac K{N(\eta)}\right)_+.
\end{equation}
For $0\le\varepsilon\le1/4$ there is a deterministic statistic of
width $C(2\sqrt\varepsilon)$ with $e(Q)\le\varepsilon$. Consequently,
writing $K(\varepsilon)=\min\{K:e_K\le\varepsilon\}$,
\begin{equation}\label{eq:v15-task-width}
 N(\eta)\left(1-\frac{2d\varepsilon}{\eta^2}\right)
       \le K(\varepsilon)\le C(2\sqrt\varepsilon).
\end{equation}
At zero error the exact width is $|v(\mathcal H)|$. It is at most the
universal marked continuation width, and it equals that width when
the event family separates all distinct marked continuation laws.
The lower bound also holds when an independent exposed selector is
available and each selector fiber has at most $K$ private labels;
in that case the conditional variance is taken given $(U,M)$.
\end{theorem}
\begin{proof}
Choose $n=N(\eta)$ separated histories and prepare them uniformly.
For an output $m$, write $\pi_h(m)=\Pp(H=h\mid M=m)$. The variance
identity gives
\begin{align*}
 \sum_{j=1}^d\E\Var(v_j(H)\mid M)
 &=\frac12\sum_m\Pp(M=m)
        \sum_{h,h'}\pi_h(m)\pi_{h'}(m)\|v(h)-v(h')\|_2^2\\
 &\ge\frac{\eta^2}{2}
       \left(1-\sum_m\Pp(M=m)\sum_h\pi_h(m)^2\right).
\end{align*}
Furthermore
\[
 \sum_m\Pp(M=m)\sum_h\pi_h(m)^2
 \le\sum_m\Pp(M=m)\max_h\pi_h(m)
 =\frac1n\sum_m\max_h Q(m\mid h)\le\frac K n.
\]
At least one coordinate variance is the average of these $d$ terms
or larger. Nonnegativity proves \eqref{eq:v15-task-packing}.
Conditioning on an independent exposed selector gives the identical
bound on the sum of coordinate variances at every fiber, with the same
uniform preparation. Average this sum over the selector and only then
choose a coordinate at least as large as its average. Thus the task
need not depend on the selector, and the same lower bound follows.

For the upper bound assign each history to a covering set containing
its feature vector and retain its set index. Every conditional
coordinate range then has length at most $2\sqrt\varepsilon$.
A real random variable in $[a,b]$ has variance at most $(b-a)^2/4$:
$\E(X-a)(b-X)\ge0$ implies
$\Var X\le(\E X-a)(b-\E X)\le(b-a)^2/4$.
This proves the upper bound uniformly over preparations.

If two histories with different feature vectors share a retained
symbol with positive probability, a preparation on just those two
histories has a strictly positive posterior variance at that symbol
in a differing coordinate. Therefore zero excess forces disjoint
retained supports for distinct feature vectors. Conversely the feature
vector itself is an exact statistic. This proves the zero law.
Every feature is determined by the universal marked continuation law;
a finite separating event family distinguishes all its finitely many
classes. The final assertion follows.
\end{proof}

\begin{corollary}[A positive-error obstruction for marked simulation]
\label{cor:v15-task-deficiency}
Suppose a prescribed source-consuming reverse simulator of width $K$
compares the stored-state restart to the full-history restart, preserving
the actual marked continuation experiment with error at most $\delta$
uniformly over preparations and announced continuation tasks. Then
\[
 \delta\ge e_K\ge
       \frac{\eta^2}{2d}\left(1-\frac K{N(\eta)}\right)_+.
\]
\end{corollary}
\begin{proof}
Compose the reverse simulator with the full-history optimal Brier
forecast for each announced task. Squared error belongs to $[0,1]$;
marked comparison changes its expectation by at most $\delta$.
The optimal stored-state forecast is no worse than this composed
forecast. Hence its excess is at most $\delta$ for every preparation
and task. Apply Theorem~\ref{thm:v15-task}.
\end{proof}

For several checkpoints these are necessary width bounds at every
restart. Their simultaneous sufficiency additionally requires that
the chosen covers admit causal updates on the reached support, exactly
as in the preceding endogenous theorem. A feature cover alone does
not construct a reverse channel preserving an arbitrary actual mark.
For a full separating family the zero-error envelope recovers the
universal continuation quotient; before completion, the quantitative
bounds depend on the declared tasks rather than on irrelevant marks
or unused future controls.
